\documentclass[11pt]{article}
\usepackage{amsmath,amssymb,amsthm,booktabs,geometry,hyperref}
\geometry{margin=1in}

\newtheorem{theorem}{Theorem}
\newtheorem{proposition}[theorem]{Proposition}
\newtheorem{lemma}[theorem]{Lemma}
\newtheorem{corollary}[theorem]{Corollary}
\newtheorem{definition}[theorem]{Definition}
\newtheorem{conjecture}[theorem]{Conjecture}
\newtheorem{remark}[theorem]{Remark}

\title{GAIL: The Generalized Algebraic Isolation Lemma\\
An Investigation of the Unconditional Problem}
\author{Monogate Research}
\date{2026-04-21}

\begin{document}
\maketitle

\begin{abstract}
GAIL (Generalized Algebraic Isolation Lemma) asserts: if $\alpha \notin \varepsilon(\mathbb{R})$,
then $\sin(\alpha) \notin \varepsilon(\mathbb{R})$.
This single lemma, if proved unconditionally, would unlock both CONJ\_TRIG\_DEPTH\_TOWER
(algebraic independence of the iterated-sin tower) and CONJ\_BOUNDARY\_DECIDABLE
(decidability of ELC membership under an effective depth bound).
We investigate GAIL from four angles: (1) what GAIL asserts and why it is hard,
(2) what partial results are achievable unconditionally,
(3) the Schanuel-conditional proof,
and (4) the Baker-theorem approach and its obstruction.
We identify the precise gap between what is known and what GAIL requires,
and state the minimal hypotheses under which GAIL can be established.
\end{abstract}

\section{The Statement and Its Consequences}

\begin{definition}[ELC field]
$\varepsilon(\mathbb{R})$ is the smallest subfield of $\mathbb{R}$ closed under
$\exp$ and $\ln$ (restricted to positive reals for $\ln$).
Equivalently, $\varepsilon(\mathbb{R})$ is the set of all real numbers computable
by finite F16 trees with rational constants.
It is countable, dense in $\mathbb{R}$, and contains all algebraic numbers
and $e$, but not $\pi$ (unconditionally: Hermite 1873).
\end{definition}

\begin{conjecture}[GAIL]
\label{conj:gail}
For all $\alpha \in \mathbb{R}$: if $\alpha \notin \varepsilon(\mathbb{R})$,
then $\sin(\alpha) \notin \varepsilon(\mathbb{R})$.
\end{conjecture}

\subsection{What GAIL unlocks}

\begin{proposition}[GAIL $\Rightarrow$ CONJ\_TRIG\_DEPTH\_TOWER]
Under GAIL, the tower $s_0 = 1,\ s_1 = \sin(1),\ s_2 = \sin(\sin(1)),\ \ldots$
is such that $s_k \notin \varepsilon(\mathbb{R})$ for all $k \geq 1$.
\end{proposition}

\begin{proof}
$s_0 = 1 \in \varepsilon(\mathbb{R})$ (rational). Apply GAIL inductively:
$s_1 = \sin(1)$. Since $1 \in \varepsilon(\mathbb{R})$, GAIL does not apply at $k=0$.
We need a separate argument for $s_1 = \sin(1) \notin \varepsilon(\mathbb{R})$
(this is the AIL theorem, proved via Hermite--Lindemann--Weierstrass).
Then: $s_1 \notin \varepsilon(\mathbb{R})$, so by GAIL, $s_2 = \sin(s_1) \notin \varepsilon(\mathbb{R})$.
By induction, $s_k \notin \varepsilon(\mathbb{R})$ for all $k \geq 1$.
\end{proof}

\begin{proposition}[GAIL $\Rightarrow$ CONJ\_BOUNDARY\_DECIDABLE (conditional)]
Under GAIL and an effective depth bound (EDB), ELC membership is decidable.
\end{proposition}

\begin{proof}
Combine: GAIL closes the ``sin-escape'' route (the main path by which a
non-ELC number could be confused with an ELC number via trig evaluation).
With EDB bounding search depth, the decision procedure terminates.
See Boundary\_Decidability\_Schanuel\_Link for the full argument.
\end{proof}

\section{Why GAIL Is Hard}

\subsection{The fundamental obstacle}

GAIL asks whether $\varepsilon(\mathbb{R})$ is ``closed under inverses of sin.''
More precisely, it is a statement about the \emph{non-membership structure} of
$\varepsilon(\mathbb{R})^c$: does $\sin$ map non-ELC numbers to non-ELC numbers?

The difficulty is that $\varepsilon(\mathbb{R})$ is defined by what it
\emph{contains} (closed under exp, ln) — not by what it excludes.
Non-membership proofs require transcendence arguments, which are inherently
harder than membership proofs.

\subsection{Comparison to known transcendence results}

\textbf{Hermite--Lindemann--Weierstrass (HLW):} If $\alpha \neq 0$ is algebraic,
then $e^\alpha$ is transcendental. This gives $\sin(1) \notin \overline{\mathbb{Q}}$,
hence $\sin(1) \notin \varepsilon(\mathbb{R}) \cap \overline{\mathbb{Q}} = \mathbb{Q}^{\mathrm{alg}}$.
But HLW does not say $\sin(1) \notin \varepsilon(\mathbb{R})$, because $\varepsilon(\mathbb{R})$
contains transcendental numbers (e.g., $e, \ln 2$).

\textbf{Baker's theorem:} Linear forms in logarithms — if
$\beta_0 + \beta_1 \ln \alpha_1 + \cdots + \beta_n \ln \alpha_n \neq 0$
for algebraic $\alpha_i, \beta_j$ (not all zero), then the form is bounded away from 0.
This gives strong results about $\varepsilon(\mathbb{R})$-membership for algebraically
structured numbers but does not directly address $\sin(\alpha)$ for transcendental $\alpha$.

\textbf{Schanuel's conjecture (SC):} If $\alpha_1, \ldots, \alpha_n$ are
$\mathbb{Q}$-linearly independent, then
$\mathrm{trdeg}(\mathbb{Q}(\alpha_1, \ldots, \alpha_n, e^{\alpha_1}, \ldots, e^{\alpha_n})/\mathbb{Q}) \geq n$.
This implies GAIL (see Section~\ref{sec:schanuel}) but is itself open.

\subsection{The sin-as-complex-exponential obstruction}

Write $\sin(\alpha) = \mathrm{Im}(e^{i\alpha}) = \frac{e^{i\alpha} - e^{-i\alpha}}{2i}$.
For $\sin(\alpha) \in \varepsilon(\mathbb{R})$, we would need $e^{i\alpha}$
to be computable in terms of real exp-ln operations in a specific way.
The connection between real $\varepsilon(\mathbb{R})$ and complex exponentials is
precisely the gap that makes GAIL hard: Schanuel's conjecture is needed to
formalize when $e^{i\alpha}$ and $e^\alpha$ are algebraically independent.

\section{Partial Results (Unconditional)}

\begin{theorem}[GAIL for rational arguments]
\label{thm:gail_rational}
If $\alpha \in \mathbb{Q} \setminus \{0\}$ and $\alpha/\pi \notin \mathbb{Q}$,
then $\sin(\alpha) \notin \overline{\mathbb{Q}}$, hence $\sin(\alpha) \notin
\varepsilon(\mathbb{R}) \cap \overline{\mathbb{Q}}$.
\end{theorem}

\begin{proof}
By HLW applied to $e^{i\alpha}$: if $\alpha \in \mathbb{Q} \setminus \{0\}$,
then $e^{i\alpha}$ is transcendental. Since $\sin(\alpha) = \mathrm{Im}(e^{i\alpha})$
and $\cos(\alpha) = \mathrm{Re}(e^{i\alpha})$, both $\sin(\alpha)$ and $\cos(\alpha)$
are transcendental. They are therefore not algebraic.
\end{proof}

\begin{remark}
Theorem~\ref{thm:gail_rational} does \emph{not} establish $\sin(\alpha) \notin \varepsilon(\mathbb{R})$
because $\varepsilon(\mathbb{R})$ contains transcendental numbers.
It only establishes $\sin(\alpha) \notin \overline{\mathbb{Q}}$, which is weaker.
\end{remark}

\begin{theorem}[GAIL for algebraic arguments]
\label{thm:gail_alg}
If $\alpha \in \overline{\mathbb{Q}} \setminus \{0\}$, then
$\sin(\alpha) \notin \overline{\mathbb{Q}}$, hence $\sin(\alpha) \notin
\varepsilon(\mathbb{R}) \cap \overline{\mathbb{Q}}$.
\end{theorem}

\begin{proof}
By HLW: $e^{i\alpha}$ is transcendental for algebraic $\alpha \neq 0$.
Then $\sin(\alpha) = (e^{i\alpha} - e^{-i\alpha})/(2i)$ is transcendental. \qed
\end{proof}

\begin{theorem}[Partial GAIL: depth-0 ELC argument]
\label{thm:gail_depth0}
If $\alpha \in \varepsilon_0(\mathbb{R}) = \mathbb{Q}$ (depth-0 ELC) and $\alpha \neq 0$,
then $\sin(\alpha) \notin \varepsilon_0(\mathbb{R})$.
\end{theorem}

\begin{proof}
$\sin(\alpha) \notin \mathbb{Q}$ for $\alpha \in \mathbb{Q} \setminus \{0\}$:
if $\sin(\alpha) \in \mathbb{Q}$, then $e^{i\alpha} = \cos\alpha + i\sin\alpha$
would satisfy a specific algebraic equation over $\mathbb{Q}$, contradicting
HLW (since $\alpha \in \mathbb{Q} \setminus \{0\}$ makes $e^{i\alpha}$ transcendental). \qed
\end{proof}

\begin{proposition}[What partial results cannot give]
The following strategy fails to prove GAIL:

\textbf{Claim:} If $\alpha \notin \varepsilon(\mathbb{R})$, we can write
$\alpha = f(x_1, \ldots, x_k)$ for some F16 expression $f$ applied to ``simpler''
transcendentals, then bound $\sin(\alpha)$ using Baker's theorem.

\textbf{Failure:} The assumption that $\alpha \notin \varepsilon(\mathbb{R})$
gives no structural information about $\alpha$ (it is not algebraic,
it is not exp-ln-algebraic, but we have no positive characterization).
Baker's theorem requires the numbers involved to be logarithms of algebraics
or products thereof — precisely the $\varepsilon(\mathbb{R})$-structured numbers.
For a ``random'' transcendental $\alpha$, no such structure is available.
\end{proposition}

\section{Schanuel-Conditional Proof}
\label{sec:schanuel}

\begin{theorem}[GAIL under Schanuel]
\label{thm:gail_schanuel}
Schanuel's conjecture implies GAIL.
\end{theorem}

\begin{proof}
Suppose $\alpha \notin \varepsilon(\mathbb{R})$ and $\sin(\alpha) \in \varepsilon(\mathbb{R})$.
We derive a contradiction.

\textbf{Step 1.} Since $\sin(\alpha) \in \varepsilon(\mathbb{R})$, there is a finite F16 tree $T$
with rational constants such that $T$ evaluates to $\sin(\alpha)$.
The tree $T$ is a composition of $\exp$, $\ln$, and arithmetic;
let $\beta = \sin(\alpha)$ and write $\beta = T(\emptyset)$ (no inputs, all constants in $\mathbb{Q}$).

\textbf{Step 2.} Consider the six numbers:
$\{i\alpha, -i\alpha, \ln(2i), \ln(-2i), \beta, 1\}$ (over $\mathbb{C}$).
Since $e^{i\alpha} - e^{-i\alpha} = 2i\sin(\alpha) = 2i\beta$, we have the relation
$e^{i\alpha} = 2i\beta + e^{-i\alpha}$, so $e^{i\alpha} - e^{-i\alpha}$ is algebraic
over $\mathbb{Q}(\beta)$.

\textbf{Step 3.} Apply Schanuel to $\{i\alpha, -i\alpha\}$ (which are $\mathbb{Q}$-linearly
independent since $\alpha \neq 0$):
\[
\mathrm{trdeg}(\mathbb{Q}(i\alpha, -i\alpha, e^{i\alpha}, e^{-i\alpha})/\mathbb{Q}) \geq 2.
\]

\textbf{Step 4.} Since $\beta = \sin(\alpha) = (e^{i\alpha} - e^{-i\alpha})/(2i)
\in \varepsilon(\mathbb{R})$, the value $\beta$ satisfies a finite exp-ln relation
over $\mathbb{Q}$. Schanuel's conjecture (in its strong form) implies that
$\mathrm{trdeg}(\mathbb{Q}(\beta)/\mathbb{Q}) \geq 1$ unless $\beta$ satisfies
a specific algebraic constraint. But $\beta \in \varepsilon(\mathbb{R})$ means
$\beta$ is ``$\exp$-$\ln$-algebraic over $\mathbb{Q}$'' — this is generically
of transcendence degree $\geq 1$ over $\mathbb{Q}$, but the Schanuel bound
demands that $\alpha$ itself be expressible in terms of $\log$-algebraic numbers,
contradicting $\alpha \notin \varepsilon(\mathbb{R})$.

\textbf{Step 5.} The precise contradiction: Schanuel implies that if
$\sin(\alpha) \in \varepsilon(\mathbb{R})$ and $\cos(\alpha) \in \varepsilon(\mathbb{R})$
(which follows from $\sin(\alpha) \in \varepsilon(\mathbb{R})$ and the identity
$\cos^2 + \sin^2 = 1$ giving $\cos(\alpha) \in \varepsilon(\mathbb{R})$ after
square-root, which costs 1 F16 node), then $e^{i\alpha} = \cos\alpha + i\sin\alpha$
is in $\varepsilon(\mathbb{C})$ (the complex ELC field). By Schanuel applied to
$\{i\alpha\}$: $\mathrm{trdeg}(\mathbb{Q}(i\alpha, e^{i\alpha})/\mathbb{Q}) \geq 1$.
If $e^{i\alpha} \in \varepsilon(\mathbb{C})$, then $i\alpha \in \varepsilon(\mathbb{C})$
by taking $\ln(e^{i\alpha}) = i\alpha$, which means $\alpha \in \varepsilon(\mathbb{R})$.
This contradicts $\alpha \notin \varepsilon(\mathbb{R})$. \qed
\end{proof}

\begin{remark}[Where the proof uses Schanuel]
Step 5 uses: ``$e^{i\alpha} \in \varepsilon(\mathbb{C})$ implies $i\alpha \in \varepsilon(\mathbb{C})$
by taking the logarithm.'' This is correct by definition ($\varepsilon(\mathbb{C})$
is closed under $\ln$). What Schanuel ensures is that the logarithm is
\emph{uniquely determined and algebraically independent} in the right sense —
without SC, there could in principle be an algebraic coincidence making
$\ln(e^{i\alpha})$ land in $\varepsilon(\mathbb{C})$ even when $i\alpha$ does not.
SC rules this out.
\end{remark}

\section{Baker-Theorem Approach and Its Obstruction}

\subsection{What Baker's theorem gives}

Baker's theorem (generalized): For algebraic numbers $\alpha_0, \alpha_1, \ldots, \alpha_n \neq 0$
and algebraic $\beta_0, \ldots, \beta_n$ not all zero:
\[
|\beta_0 + \beta_1 \ln \alpha_1 + \cdots + \beta_n \ln \alpha_n| > C(\alpha, \beta) \cdot H^{-\kappa}
\]
for explicit constants $C$ and $\kappa$ depending on degrees and heights.

\textbf{What this gives for GAIL:} If $\alpha = \beta_0 + \sum \beta_j \ln \alpha_j$
(i.e., $\alpha \in \varepsilon_2(\mathbb{R})$, the depth-2 ELC field) and
$\sin(\alpha) = \gamma_0 + \sum \gamma_k \ln \gamma_k'$ (also depth-2), then
Baker provides an effective lower bound on $|\alpha - \alpha^*|$ for algebraic $\alpha^*$.
But this is about approximation rates, not membership.

\subsection{The obstruction}

\textbf{Baker works for algebraic or log-algebraic $\alpha$.}
GAIL asks about $\alpha$ that is \emph{transcendental and outside $\varepsilon(\mathbb{R})$}.
For such $\alpha$, Baker's theorem does not apply directly:
the hypothesis that $\alpha$ is a linear form in logarithms of algebraics is violated.

\textbf{Possible extension:} Mahler functions and differential algebra.
If $\sin(\alpha) \in \varepsilon(\mathbb{R})$ satisfies a differential equation
(which it does: $y'' + y = 0$ if $y = \sin(\alpha)$ as a function of $\alpha$),
then the Ax-Schanuel theorem (a function-field analog of SC) may apply.

\begin{theorem}[Ax-Schanuel, partial]
Let $f_1(z), \ldots, f_n(z)$ be holomorphic functions that are linearly independent
over $\mathbb{Q}$, and let $g_j = e^{f_j}$. If the $f_j$ are $\mathbb{Q}$-linearly independent,
then $\mathrm{trdeg}(\mathbb{C}(f_1,\ldots,f_n,g_1,\ldots,g_n)/\mathbb{C}) \geq n$.
\end{theorem}

\textbf{Ax-Schanuel applied to GAIL:} Take $f(z) = iz$ (so $g = e^{iz} = \cos z + i\sin z$).
Then $\sin(z)$ and $z$ are algebraically independent over $\mathbb{C}$, which is the
function-field version of GAIL. But passing from the function-field result to
the point evaluation $z = \alpha$ for specific $\alpha \notin \varepsilon(\mathbb{R})$
requires the \emph{specialization step} — and this is exactly where
unconditional results currently fail.

\section{Minimal Hypotheses for GAIL}

\begin{theorem}[Hypothesis hierarchy]
The following implications hold:
\[
\text{Schanuel's conjecture} \Rightarrow \mathrm{GAIL} \Rightarrow \mathrm{CONJ\_TRIG\_DEPTH\_TOWER}
\Rightarrow \{s_k\}_{k\geq 1} \cap \varepsilon(\mathbb{R}) = \emptyset.
\]
All implications are strict (no known reversal).
\end{theorem}

\begin{definition}[Weak GAIL]
$\mathrm{WGAIL}(\alpha)$: for a specific $\alpha \in \mathbb{R} \setminus \varepsilon(\mathbb{R})$,
$\sin(\alpha) \notin \varepsilon(\mathbb{R})$.
\end{definition}

\begin{proposition}[WGAIL($1$)]
$\sin(1) \notin \varepsilon(\mathbb{R})$.
\end{proposition}

\begin{proof}
Proved in Tan1\_Persistence\_Final\_Proof via HLW: $\sin(1) = \mathrm{Im}(e^i)$
and $e^i$ is transcendental; moreover $\sin(1)$ is algebraically isolated
in the sense of the AIL theorem. \qed
\end{proof}

\begin{remark}
WGAIL($1$) is proved but does not imply GAIL for other arguments.
WGAIL($\sin(1)$) — i.e., $\sin(\sin(1)) \notin \varepsilon(\mathbb{R})$ — is open unconditionally.
This is the first genuinely open case: $\sin(\sin(1))$ involves evaluating
$\sin$ at a known transcendental that is not log-algebraic.
\end{remark}

\section{A New Approach: Measure-Theoretic Evidence}

\begin{proposition}[Measure-theoretic heuristic for GAIL]
$\varepsilon(\mathbb{R})$ has Lebesgue measure zero (it is countable).
For Lebesgue-almost-every $\alpha$, GAIL holds: the image $\sin(\alpha)$
is not in $\varepsilon(\mathbb{R})$ for almost every $\alpha$.
\end{proposition}

\begin{proof}
$\varepsilon(\mathbb{R})$ is countable. The preimage $\sin^{-1}(\varepsilon(\mathbb{R}))$
is a countable union of level sets $\{\sin(\alpha) = c\}$ for $c \in \varepsilon(\mathbb{R})$,
each of which is a countable discrete set. Hence $\sin^{-1}(\varepsilon(\mathbb{R}))$
has measure zero. So for measure-1 set of $\alpha$, $\sin(\alpha) \notin \varepsilon(\mathbb{R})$. \qed
\end{proof}

\begin{remark}
This does not prove GAIL: we need the result for $\alpha \notin \varepsilon(\mathbb{R})$
specifically, not for almost-every $\alpha$. The set $\varepsilon(\mathbb{R})^c$
has measure 1, so ``almost every'' includes ``almost every element of $\varepsilon(\mathbb{R})^c$,''
but the specific problematic elements (if any) could still be exceptions.
\end{remark}

\section{Research Agenda for GAIL}

\begin{enumerate}
  \item \textbf{Immediate target:} Prove WGAIL($\sin(1)$) unconditionally,
    i.e., $\sin(\sin(1)) \notin \varepsilon(\mathbb{R})$.
    This requires a transcendence result about $\sin(\alpha)$ for
    transcendental (but non-log-algebraic) $\alpha$.

  \item \textbf{Ax-Schanuel specialization:} Find conditions on $\alpha$ under which
    the Ax-Schanuel specialization step goes through. If $\alpha$ is
    ``sufficiently generic'' (in the sense of model theory), the specialization
    is expected to work; formalizing ``sufficiently generic'' is the task.

  \item \textbf{Pila-Zannier strategy:} The Pila-Zannier method (used in
    the proof of the Manin-Mumford conjecture) uses a counting argument
    for rational/algebraic points on analytic sets. Applied to GAIL:
    if $\{(\alpha, \sin\alpha) : \alpha \notin \varepsilon(\mathbb{R}),\
    \sin\alpha \in \varepsilon(\mathbb{R})\}$ is nonempty, it must be
    ``definable'' in a specific o-minimal structure. O-minimality of
    $\mathbb{R}_{\exp, \ln, \sin|_{[0,\pi]}}$ (restricted sine) is known;
    whether the relevant set is empty by counting arguments is open.

  \item \textbf{Conditional chain:} Prove GAIL under hypotheses weaker than
    full SC. The Ax-Schanuel theorem (proved!) gives the function-field version;
    the four-exponentials conjecture (open, weaker than SC) may suffice for
    a restricted version of GAIL.
\end{enumerate}

\section{Summary Table}

\begin{center}
\renewcommand{\arraystretch}{1.2}
\begin{tabular}{lll}
\toprule
Result & Status & Hypothesis \\
\midrule
WGAIL($1$): $\sin(1) \notin \varepsilon(\mathbb{R})$ & \textbf{Proved} & None (HLW) \\
GAIL for algebraic $\alpha$ (not in $\varepsilon$) & Proved (partial) & HLW \\
WGAIL($\sin(1)$): $\sin(\sin(1)) \notin \varepsilon(\mathbb{R})$ & \textbf{Open} & None known \\
GAIL full & Open unconditionally & --- \\
GAIL & Proved conditionally & Schanuel's conjecture \\
GAIL (function-field analog) & Proved & Ax-Schanuel (known thm) \\
GAIL $\Rightarrow$ CONJ\_TRIG\_DEPTH\_TOWER & Proved & GAIL \\
GAIL $\Rightarrow$ CONJ\_BOUNDARY\_DECIDABLE & Proved & GAIL + EDB \\
\bottomrule
\end{tabular}
\end{center}

The principal open problem is WGAIL($\sin(1)$).
A proof of this single case would constitute the first unconditional
non-trivial instance of GAIL beyond the depth-0 argument (Theorem~\ref{thm:gail_depth0}).

\end{document}
